\section{Introduction}
Citizens increasingly ask general-purpose LLMs questions of law. The models
answer in fluent, authoritative prose---and sometimes cite articles that do
not exist, or that exist but say something else. Unlike open-domain chit-chat,
the legal domain has an asymmetric cost structure: a plausible-sounding but
fabricated citation actively \emph{misleads}, and is worse than an admission
of ignorance. Hallucination in natural language generation is by now
well documented~\cite{ji2023hallucination}, and retrieval augmentation is the
standard mitigation~\cite{lewis2020rag}; yet grounding a generation in
retrieved passages \emph{reduces} but does not \emph{guarantee} the absence of
invented references.

\emph{Le Rapporteur} takes a deliberately conservative stance suited to law.
The system is organized around a single invariant---\textbf{no verifiable
source, no answer}---enforced not by prompt engineering but by a
post-generation verification stage that is \emph{fail-closed}: any doubt,
network error, or unresolved reference collapses to an explicit refusal
(``Je ne trouve pas de texte applicable.'', \emph{I find no applicable
text.}). The design is a defense in depth: the model may draft freely, but
nothing reaches the user unless each citation resolves to a real, in-force
article in an authoritative database.

\paragraph{Contributions.}
\begin{itemize}
  \item A \textbf{fail-closed pipeline} for legal QA (\emph{generate,
  extract, verify, source-or-refuse}) whose trust guarantee is a property of
  the \emph{system}, not of the model (\S\ref{sec:overview}).
  \item A \textbf{database-centric citation verifier} over the Canutes/Légifrance
  corpus, and the query engineering (index-friendly citation normalization,
  state-aware selection of the in-force version) that makes per-citation
  checking feasible at interactive latency over a large, irregular schema
  without full-table regular-expression scans (\S\ref{sec:verify}).
  \item A \textbf{reusable trust service}: the verifier is exposed over HTTP
  (\texttt{/verify}) and as a Model Context Protocol (MCP) server, so any
  third-party agent can fact-check the citations in an arbitrary LLM output
  (\S\ref{sec:service}).
  \item A \textbf{hardware-sovereign serving layer}: an OpenAI-compatible
  indirection lets the identical pipeline run on non-NVIDIA accelerators; we
  demonstrate it live on a Qualcomm Cloud~AI~100 and target AMD via a portable
  runtime (\S\ref{sec:sovereign}).
  \item An \textbf{executable trust specification} in Gherkin, readable and
  auditable by a jurist, in which a single fabricated article fails the build
  (\S\ref{sec:bdd}).
\end{itemize}

We are candid throughout about what is \emph{proven} versus \emph{envisioned};
this project was built under hackathon constraints and we mark roadmap items
explicitly.
